\subsection{Lokale Diffeomorphismen}

\begin{definition}{}
Seien $ U, V \subset \mathbb{R}^n $ offene Mengen. Eine Abbildung $ \varphi: U \to V $ heißt \emph{Diffeomorphismus}, falls $ \varphi $ bijektiv, stetig differenzierbar ist und die Umkehrfunktion $ \varphi^{-1}: V \to U $ ebenfalls stetig differenzierbar ist.
\end{definition}

\begin{bemerkungbox}
Falls $ \varphi $ ein Diffeomorphismus ist, gilt für die Jacobi-Matrix


$$
J(\varphi^{-1}(x)) = (J\varphi(x))^{-1}.
$$


Insbesondere folgt daraus, dass die Inverse einer $ C^1 $-Funktion ebenfalls $ C^1 $ ist.
\end{bemerkungbox}

\begin{definition}{}
Eine Funktion $ \varphi: D \to \mathbb{R}^n $, wobei $ D \subset \mathbb{R}^n $ offen ist, heißt \emph{lokaler Diffeomorphismus}, falls jeder Punkt $ x \in D $ eine Umgebung $ U \subset D $ besitzt, sodass $ \varphi $ eingeschränkt auf $ U $ ein Diffeomorphismus ist.
\end{definition}

\begin{satz}{}
Sei $ D \subset \mathbb{R}^n $ offen und sei $ \varphi: D \to \mathbb{R}^n $ eine $ C^1 $-Funktion. Falls die Jacobi-Matrix $ J\varphi(a) $ in einem Punkt $ a \in D $ invertierbar ist, das heißt


$$
\det J\varphi(a) \neq 0,
$$


so existiert eine Umgebung $ U \subset D $ von $ a $, sodass $ \varphi $ eingeschränkt auf $ U $ ein Diffeomorphismus ist.
\end{satz}

\begin{beweisbox}
Die Idee des Beweises beruht darauf, dass eine lokal invertierbare Funktion eine lokal definierte Umkehrfunktion besitzt. Mithilfe der Stetigkeit von $ J\varphi(x) $ kann gezeigt werden, dass $ \varphi $ in der Umgebung von $ a $ eine eindeutige Umkehrfunktion besitzt.
\end{beweisbox}


% Spätere Vorlesung

\begin{satz}{}
Es sei $ D \subset \mathbb{R}^n $ offen und $ f: D \to \mathbb{R}^n $ eine $ C^1 $-Funktion. Falls die Jacobi-Matrix $ Jf(a) $ in einem Punkt $ a \in D $ invertierbar ist, das heißt


$$
\det Jf(a) \neq 0,
$$


so existiert eine Umgebung $ U \subset D $ von $ a $, sodass $ f $ eingeschränkt auf $ U $ ein Diffeomorphismus ist.
\end{satz}

\begin{beweisbox}
Die Beweisidee basiert auf dem Fixpunktsatz von Banach. Es wird gezeigt, dass eine lokal invertierbare Funktion eine lokal definierte Umkehrfunktion besitzt. Durch die Stetigkeit von $ Jf(x) $ kann die Existenz einer eindeutigen Lösung gesichert werden.
\end{beweisbox}

\begin{motivationbox}
Die Taylor-Entwicklung ermöglicht die lokale Approximation von Funktionen. Falls $ Jf(a) $ invertierbar ist, kann die Funktion um $ a $ gut linear approximiert werden, was die Voraussetzung für die Anwendung des Fixpunktsatzes erfüllt.
\end{motivationbox}

\begin{satz}{}
Sei $ \varphi $ ein kontrahierender Operator auf einem vollständigen metrischen Raum. Dann besitzt $ \varphi $ einen eindeutigen Fixpunkt $ x^* $ mit $ \varphi(x^*) = x^* $.
\end{satz}

\begin{beweisbox}
Der Beweis verwendet eine iterative Konstruktion. Für eine Startfolge $ x_n $ zeigt man, dass diese eine Cauchy-Folge bildet und ihr Grenzwert die Fixpunktbedingung erfüllt.
\end{beweisbox}

\begin{satz}{}
Sei $ f $ eine differenzierbare Funktion mit einer invertierbaren Jacobi-Matrix. Dann ist $ f $ lokal umkehrbar und die Umkehrfunktion ist ebenfalls differenzierbar.
\end{satz}

\begin{beweisbox}
Die Beweisidee nutzt die Abschätzung der Ableitungen von $ f $ und deren Invertierbarkeit, um die Existenz einer lokalen Umkehrfunktion nachzuweisen. Die Ableitungsabschätzungen zeigen zudem die Differenzierbarkeit der Umkehrfunktion.
\end{beweisbox}
